\section{The integral conclusion and motivic consequences}\label{sec:conclusion}

\subsection{The integral zero-cycle theorem}\label{subsec:integral-proof}

Both numerical branches now supply the nonzero coefficient required
by Corollary~\ref{rel:correspondence}. We assemble the rational
vanishing and then use Roitman's theorem to remove torsion.

\begin{proof}[Proof of Theorem~\ref{thm:bloch}]
For a surface of Kodaira dimension below two, the assertion follows
from Bloch--Kas--Lieberman as in
Proposition~\ref{prop:reduction}. For a surface of general type,
pass to a smooth minimal model \(X_0\). If \(K_{X_0}^2\le8\),
Proposition~\ref{prop:nonrank} gives a nonzero diagonal coefficient
in Theorem~\ref{thm:chow-relation} on a point blowup of \(X_0\).
If \(K_{X_0}^2=9\), Proposition~\ref{prop:rankone} gives such a
coefficient on \(X_0\) itself. In either case
Corollary~\ref{rel:correspondence} and Lemma~\ref{lem:birational} yield
\[
\CH_0(S)^0\otimes_{\mathbb Z}\mathbb Q=0.
\]

Let \(z\in\CH_0(S)^0\). Vanishing after tensoring with
\(\mathbb Q\) means that some nonzero integer annihilates \(z\).
Roitman's torsion Theorem makes the Albanese map injective on torsion
zero-cycles over \(\mathbb C\)
\cite[Theorem~4.2]{BlochRoitman}; see also \cite{Roitman}.
But \(\Alb(S)=0\), since \(q(S)=0\). Therefore \(z=0\).
Thus the integral degree map is injective. Any complex point of
\(S\) has degree one, so it is also surjective.
\end{proof}

\begin{proof}[Proof of Corollary~\ref{cor:bloch-pg-zero}]
If \(S\) is not of general type, Bloch, Kas, and Lieberman prove that
the kernel of \(\operatorname{alb}_S\) vanishes with integral
coefficients \cite[pp.~135--136]{BKL}. Surjectivity gives the claim.

Suppose \(S\) is of general type, and let \(S_0\) be a smooth minimal
model. The inequalities \(K_{S_0}^2>0\) and \(\deg c_2(S_0)>0\),
together with Noether's formula, give
\(1-q(S_0)+p_g(S_0)>0\) \cite[Remark~1.3.1]{PedriniWeibel}.
By Lemma~\ref{lem:birational}, \(p_g(S_0)=p_g(S)=0\) and
\(q(S_0)=q(S)\), so the inequality forces \(q(S)=0\).
Theorem~\ref{thm:bloch} now gives \(\CH_0(S)^0=0\), while
\(\Alb(S)(\mathbb C)=0\). Thus \(\operatorname{alb}_S\) is an
isomorphism.
\end{proof}

\subsection{Rational Chow motives}\label{subsec:motivic-consequences}

We now pass from integral zero-cycles to rational Chow motives.
Corollary~\ref{cor:bloch-pg-zero} supplies the Albanese-kernel
vanishing for every surface with \(p_g=0\); the additional hypothesis
\(q=0\) enters only in the pure-Tate conclusions below.

\begin{proof}[Proof of Corollary~\ref{cor:tate-motives}]
Let \(T(S)\) be the kernel of the rational Albanese map
\[
\CH_0(S)^0\otimes_{\mathbb Z}\mathbb Q
\longrightarrow \Alb(S)(\mathbb C)\otimes_{\mathbb Z}\mathbb Q.
\]
Corollary~\ref{cor:bloch-pg-zero} gives \(T(S)=0\).
Kahn--Murre--Pedrini's equivalence
\cite[Corollaries~7.4.9 and~7.6.11]{KMP} then gives \(t_2(S)=0\)
and Kimura finite-dimensionality of \(h(S)_{\mathbb Q}\). Their
hypotheses hold over \(\mathbb C\), an algebraically closed
characteristic-zero field of infinite transcendence degree over
\(\mathbb Q\). This step does not require \(q(S)=0\).

Now assume \(q(S)=0\). Kahn--Murre--Pedrini's surface
Chow--K\"unneth decomposition
\cite[Propositions~7.2.1 and~7.2.3]{KMP} is
\[
h(S)_{\mathbb Q}\simeq
\mathbf 1\oplus h_1(S)\oplus\mathbb L^{\oplus\rho(S)}
\oplus t_2(S)\oplus h_3(S)\oplus\mathbb L^2,
\]
where \(\rho(S)\) is the Picard number. The endomorphism
identifications for \(h_1(S)\) and \(h_3(S)\) with those of the
Albanese and Picard varieties show that these two motives vanish when
\(q(S)=0\). The Betti realization of the middle decomposition, with
\(t_2(S)=0\), gives \(\rho(S)=b_2(S)\). This gives the asserted
decomposition. Tensoring it proves the assertion for all powers of
\(S\).

For a partition \(\nu=1^{a_1}\cdots n^{a_n}\) of \(n\ge1\), put
\(\ell(\nu)=\sum_j a_j\), \(\Sigma_\nu=\prod_j\mathfrak S_{a_j}\),
and \(S^{(\nu)}=S^{\ell(\nu)}/\Sigma_\nu\).
The formula of de Cataldo--Migliorini
\cite[Theorem~6.2.1]{deCataldoMigliorini} gives
\[
h(S^{[n]})_{\mathbb Q}\simeq
\bigoplus_{\nu\vdash n}h(S^{(\nu)})_{\mathbb Q}
 \otimes\mathbb L^{\,n-\ell(\nu)}.
\]
Here \(h(S^{(\nu)})_{\mathbb Q}\) is the rational quotient motive of
the possibly singular symmetric-product quotient. In the quotient
formalism of \cite[Section~4.3, Lemma~4.3.1, and Remark~5.4.5]{deCataldoMigliorini},
it is the image on \(h(S^{\ell(\nu)})_{\mathbb Q}\) of the averaging
projector
\[
e_\nu=\frac{1}{|\Sigma_\nu|}
\sum_{\sigma\in\Sigma_\nu}\Gamma_\sigma.
\]
Since \(\Hom(\mathbb L^a,\mathbb L^b)\) is zero for \(a\ne b\) and
is \(\mathbb Q\) for \(a=b\), every idempotent summand of a finite
Tate sum is again a finite Tate sum. Each term in the formula is
therefore pure Tate, as is \(h(S^{[n]})_{\mathbb Q}\). The case
\(n=0\) is immediate.

Finally, on a summand \(\mathbb L^a\), the rational Chow group and
even Betti cohomology each have one generator in codimension \(a\)
and vanish in other codimensions; cycle class identifies these
generators, and odd Betti cohomology vanishes. Compatibility of cycle
class with correspondences transfers these assertions to the stated
pure Tate motives. Thus the displayed cycle-class maps are
isomorphisms of rational vector spaces. This uses the decomposition
in rational Chow motives, not finite-dimensionality alone.
\end{proof}
